% Zumdahl Chapter 9 free-response set -- 2 long (10) + 3 short (4) = 32 pts.
\documentclass{apchem-exam}

\apcunitword{Chapter}
\apcunit{9}
\apcunittitle{Covalent Bonding: Orbitals}
\apcdoctitle{Free-Response Set}
\apcsubtitle{Zumdahl \S9.1 \textbullet\ 5 questions \textbullet\ 32 points \textbullet\ 50 minutes}

\begin{document}
\apctitleblock

\begin{apcnote}[Scope --- most of this chapter is off the framework]
Only \S9.1 is examinable, as the orbital half of \textbf{CED 2.7}. Read
the two exclusions before assigning:

\textbf{1. Hybridization involving $d$ orbitals is excluded.} For a
central atom with more than four electron domains the CED requires the
\emph{shape} only. \ce{sp^3d} and \ce{sp^3d^2} labels are not assessed and
are not asked for here.

\textbf{2. Molecular orbital theory is excluded.} \S9.2--9.5 --- MO
diagrams, bond order from MO theory, paramagnetism of \ce{O2} --- are
enrichment. They are genuinely illuminating and they will not be on the
exam. No question below uses them.

What remains is worth doing properly: hybridization of second-row central
atoms, and the sigma/pi count that explains why double bonds cannot
rotate.
\end{apcnote}

\examsection{II}{Free Response}{5 questions \textbullet\ 32 points \textbullet\ 50 minutes}{\calculatorok}

% =====================================================================
\frq{long}{10}{Zumdahl \S9.1 \textbullet\ CED 2.7 \textbullet\ hybridization from electron domains}

Consider the molecules \ce{CH4}, \ce{C2H4} (ethene) and \ce{C2H2}
(ethyne).

\begin{parts}
  \item State the number of electron domains around each carbon atom and
        the hybridization that follows.
        \answerlines[3]{\ce{CH4}: 4 domains, \ce{sp^3}. \quad
        \ce{C2H4}: 3 domains on each carbon, \ce{sp^2}. \quad
        \ce{C2H2}: 2 domains on each carbon, \ce{sp}.}
  \item State the bond angle and geometry about carbon in each.
        \answerlines[3]{\ce{CH4}: tetrahedral, \SI{109.5}{\degree}. \quad
        \ce{C2H4}: trigonal planar, \SI{120}{\degree}. \quad
        \ce{C2H2}: linear, \SI{180}{\degree}.}
  \item Explain the rule that connects the number of domains to the
        hybridization. \workspace[2.8cm]
        \keyonly{\begin{apcanswer}Each electron domain needs its own
        orbital pointing at it. The atom therefore mixes exactly as many
        atomic orbitals as it has domains: two domains mix one $s$ and one
        $p$ (\ce{sp}), three mix one $s$ and two $p$ (\ce{sp^2}), four mix
        one $s$ and three $p$ (\ce{sp^3}).

        The count of orbitals mixed always equals the count of domains,
        because orbitals are conserved --- $n$ atomic orbitals give
        exactly $n$ hybrid orbitals.\end{apcanswer}}
  \item Count the sigma and pi bonds in ethene, \ce{C2H4}.
        \answerlines[3]{\textbf{Five sigma} --- four \ce{C-H} plus one
        \ce{C-C} --- and \textbf{one pi}, the second component of the
        \ce{C=C}. A double bond is always one sigma plus one pi, never two
        sigmas.}
\end{parts}

\begin{rubric}
  \rpoint{3}{(a) 1 pt each molecule, requiring both the domain count and
    the hybrid label}
  \rpoint{2}{(b) 1 pt all three geometries; 1 pt all three angles}
  \rpoint{2}{(c) 1 pt one orbital per domain; 1 pt orbitals are conserved,
    $n$ in gives $n$ out}
  \rpoint{3}{(d) 1 pt five sigma; 1 pt one pi; 1 pt the double bond
    identified as one sigma plus one pi. Answering ``two pi'' earns 0 for
    the last two points}
\end{rubric}

% =====================================================================
\frq{long}{10}{Zumdahl \S9.1 \textbullet\ CED 2.7 \textbullet\ sigma, pi, and restricted rotation}

\begin{parts}
  \item Describe how a sigma bond and a pi bond differ in the way their
        orbitals overlap. \workspace[3cm]
        \keyonly{\begin{apcanswer}A \textbf{sigma} bond forms by
        \emph{end-on} overlap along the axis joining the two nuclei. Its
        electron density is concentrated directly between the nuclei and
        is cylindrically symmetric about that axis.

        A \textbf{pi} bond forms by \emph{side-on} overlap of two parallel
        $p$ orbitals. Its density lies in two lobes above and below the
        internuclear axis, with a nodal plane containing the
        axis.\end{apcanswer}}
  \item Explain why rotation about the \ce{C=C} bond in ethene is
        restricted, while rotation about the \ce{C-C} bond in ethane is
        essentially free. \workspace[3.2cm]
        \keyonly{\begin{apcanswer}Ethane's \ce{C-C} bond is a single sigma
        bond. Because sigma overlap is cylindrically symmetric about the
        bond axis, twisting one carbon relative to the other does not
        reduce the overlap, so rotation costs almost nothing.

        Ethene's double bond includes a pi bond formed by two parallel $p$
        orbitals. Rotating one \ce{CH2} group turns those orbitals out of
        alignment; at \SI{90}{\degree} they are perpendicular and the
        side-on overlap is zero. Rotation therefore requires breaking the
        pi bond, which is energetically expensive --- so the two ends stay
        locked in one plane.\end{apcanswer}}
  \item State one observable chemical consequence of that restricted
        rotation. \answerlines[2]{Molecules such as
        1,2-dichloroethene exist as two distinct, separable isomers (the
        two chlorines on the same side, or on opposite sides), which would
        be impossible if the double bond rotated freely.}
  \item State the hybridization of the nitrogen atom in \ce{NH3}, and
        explain why the lone pair must be counted.
        \answerlines[3]{\ce{sp^3}. The lone pair is an electron domain
        occupying an orbital of its own and taking up space around the
        nitrogen, so it must be counted when determining hybridization ---
        four domains, hence \ce{sp^3}. Counting only the three bonds would
        wrongly give \ce{sp^2}.}
\end{parts}

\begin{rubric}
  \rpoint{2}{(a) 1 pt sigma end-on along the internuclear axis; 1 pt pi
    side-on, above and below the axis}
  \rpoint{4}{(b) 2 pts sigma symmetry means overlap is unchanged by
    rotation; 2 pts rotation destroys pi overlap and would require
    breaking the pi bond. ``Double bonds are stronger'' earns at most 1}
  \rpoint{2}{(c) any correct isomerism consequence}
  \rpoint{2}{(d) 1 pt \ce{sp^3}; 1 pt the lone pair counts as a domain}
\end{rubric}

% =====================================================================
\frq{short}{4}{Zumdahl \S9.1 \textbullet\ CED 2.7 \textbullet\ assigning hybridization}

State the hybridization of the central atom in each species.

\begin{parts}
  \item \ce{BF3} \answerlines[2]{Three domains, no lone pair on boron:
        \ce{sp^2}.}
  \item \ce{H2O} \answerlines[2]{Four domains (two bonds, two lone pairs):
        \ce{sp^3}.}
  \item \ce{CO2} \answerlines[2]{Two domains on carbon: \ce{sp}.}
\end{parts}

\begin{rubric}
  \rpoint{4}{1 pt each correct hybridization, plus 1 pt for consistently
    counting lone pairs as domains across all three. A candidate who
    answers \ce{sp^2} for water has counted only the bonds}
\end{rubric}

% =====================================================================
\frq{short}{4}{Zumdahl \S9.1 \textbullet\ CED 2.7 \textbullet\ counting sigma and pi bonds}

\begin{parts}
  \item State the number of sigma and pi bonds in \ce{C2H2} (ethyne).
        \answerlines[2]{Three sigma (two \ce{C-H}, one \ce{C-C}) and
        \textbf{two} pi --- a triple bond is one sigma plus two pi.}
  \item State the number of sigma and pi bonds in \ce{HCN}.
        \answerlines[2]{Two sigma (\ce{H-C} and \ce{C-N}) and two pi, in
        the \ce{C#N} triple bond.}
\end{parts}

\begin{rubric}
  \rpoint{2}{(a) 1 pt three sigma; 1 pt two pi}
  \rpoint{2}{(b) 1 pt two sigma; 1 pt two pi. The general rule --- the
    first bond of any pair is sigma, every additional bond is pi --- earns
    full credit wherever it is applied correctly}
\end{rubric}

% =====================================================================
\frq{short}{4}{Zumdahl \S9.1 \textbullet\ CED 2.7 \textbullet\ shape with more than four domains}

Sulfur hexafluoride, \ce{SF6}, has six electron domains around sulfur.

\begin{parts}
  \item State the molecular shape and the \ce{F-S-F} bond angle.
        \answerlines[2]{\textbf{Octahedral}, with \SI{90}{\degree} angles
        between adjacent fluorines.}
  \item Explain why six domains adopt this arrangement.
        \answerlines[3]{Electron domains repel one another, so they
        arrange themselves as far apart as possible on the surface of a
        sphere around the central atom. For six identical domains that
        arrangement is octahedral, which places each domain
        \SI{90}{\degree} from its four neighbours and \SI{180}{\degree}
        from the one opposite.}
\end{parts}

\begin{rubric}
  \rpoint{2}{(a) 1 pt octahedral; 1 pt \SI{90}{\degree}}
  \rpoint{2}{(b) 1 pt domains repel and maximize separation; 1 pt the
    octahedral result described. \textbf{Do not require or credit an
    \ce{sp^3d^2} label} --- $d$-orbital hybridization is excluded by the
    CED for topic 2.7}
\end{rubric}

\examtotal

\end{document}
